\section{Introduction and historical evolution}
The concept of virtual machines was initially developed in the 1980s at the IBM laboratories as a method of partitioning expensive mainframe systems to allow the simultenous execution of multiple independent operating systems, thus maximizing the utilisation of hardware resources and reduction of mainframe costs \cite{reuben2007virtual, smith2005virtual}.The virtual machine monitor or hypervisor acts as an intermediary layer between the physical hardware platform and the guest operating systems, creating the ilusion of a dedicated physical machine for each guest instance \cite{king2003operating,smith2005virtual}.

\subsection{Popek-Goldberg theorem}
In 1974 Gerald J. Popek and Robert P. Goldberg published the paper \emph{"Formal Requirements for Virtualizable Third Generation Architectures"} \cite{popek1974formal} in which they defined the strict criterias that a processor architecture must meet to be considered natively virtualizable without modifying the guest system operating system. Acording to their theorem, a processor's instructions are divided into three categories as follow: 
    \begin{itemize}
        \item \textbf{Privileged instructions}: These are the instructions that generate a hardware trap if executed in user mod also known as unprivileged, transferring control to the operating system kernel.
        \item \textbf{Control sensitive instructions}: These are the instructions that attempt to modify of the physical resources or the system state.
        \item \textbf{Behaviour sensitive instructions}: These are the instructions whose execution depends on the current state of the hardware or that reveal the privilege level at which code is running.
    \end{itemize}

Based on these definitions, Popek-Goldberg theorem mentions that an computere architectue is virtualizable in native mode if and only if the sum of sensible instructions is a strict subsum of the sum of all the privileged instructions. If an architecture respects this theorem, any attempt by the guest system to execute a sensitive instruction will automatically generate a hardware trap in the hypervisor. The hypervisor (which runs in supervisor mode) can safely emulate the behavior of that instruction, providing the virtual machine with the illusion of complete control. The original x86 architecture famously violated this rule (having 17 sensitive but unprivileged instructions, such as POPF, PUSHF, SGDT, SLDT), which necessitated complex software techniques (dynamic binary translation) until the introduction of Intel VT-x and AMD-V hardware extensions in 2005 \cite{vm-security.pdf}.

\subsection{Virtual machines taxonomy}
In particular, virtual machines are classified according to the level of abstraction they provide and the communication interfaces they expose \cite{vm-security.pdf}. Process Virtual Machines, they are created to run a single user program. They provide an interface at the ABI (Application Binary Interface) level and isolate the process from the host operating system. Examples include the JVM (Java Virtual Machine), CLI (Common Language Infrastructure) virtual machines in the Microsoft .NET ecosystem and Docker deployments. System Virtual Machines, they provide complete emulation of a hardware platform, allowing the installation and execution of a completely independent operating system. These are further classified into Type I hypervisors, bare-metal, running directly on hardware, such as VMware ESXi or Xen and Type II, hosted, running on top of a host OS, such as Oracle VirtualBox, or VMware Workstation \cite{king.pdf}.

